%% ---------------------------------------------------------------------
\section{Por qué no hay oscilaciones en una dimensión}
\label{sec:un-monotonia}
%% ---------------------------------------------------------------------

La línea de fase deja ver algo más fuerte que la clasificación local: sobre la recta, una solución solo puede avanzar en un sentido.

\begin{proposicion}[label=prop:un-monotonia]{Monotonía de las soluciones}
Sea $f$ localmente Lipschitz (por ejemplo, $f\in\Cont1$) y sea $x(t)$ la solución de \eqref{eq:un-escalar} en su intervalo máximo de existencia $[0,T)$.
\begin{enumerate}
  \item Si $f(x_0)=0$, entonces $x(t)\equiv x_0$.
  \item Si $f(x_0)>0$ (resp. $<0$), entonces $x(t)$ es estrictamente creciente (resp. decreciente) y $f(x(t))\neq0$ para todo $t\in[0,T)$.
  \item Si además $x(t)$ es acotada, entonces $T=\infty$ y $x(t)$ converge a un equilibrio cuando $t\to\infty$.
\end{enumerate}
\end{proposicion}

\begin{proof}
(1) La constante $x_0$ es una solución y, por la unicidad (\cref{thm:pm-local}), es la única. (2) Si existiera $t_1$ con $f(x(t_1))=0$, la constante $x(t_1)$ y la solución $x(t)$ resolverían el mismo problema de valor inicial en $t_1$, lo que contradice la unicidad. Luego $f(x(t))$ no se anula y, por continuidad, conserva el signo de $f(x_0)$; como $\dot x=f(x)$, sigue la monotonía. (3) Que $T=\infty$ lo da el \cref{thm:pm-compacto}. Una función monótona y acotada tiene límite $\ell$. Si $f(\ell)\neq0$, por ejemplo $f(\ell)>0$, entonces $\dot x\ge f(\ell)/2$ para $t$ grande y $x$ crecería sin cota, lo que es absurdo. Luego $f(\ell)=0$.
\end{proof}

La consecuencia es inmediata: \alerta{un sistema escalar invariante en el tiempo, con $f$ localmente Lipschitz, no tiene soluciones periódicas no constantes}. Las oscilaciones exigen al menos dos dimensiones (capítulo~4) o un estado que viva en una circunferencia en lugar de una recta (\cref{cap:orbitas}).

\begin{trampa}[label=tr:un-oscilar]{En una dimensión nada oscila}
\plausible{Un sistema con una sola variable de estado no puede oscilar.}
La conclusión anterior depende de tres hipótesis, y falla si se quita cualquiera:
\begin{itemize}
  \item \emph{Invariancia en el tiempo.} $\dot x=-x+\sin t$ es escalar y sus soluciones tienden a la oscilación $\tfrac12(\sin t-\cos t)$.
  \item \emph{Estado en la recta.} $\dot\theta=\omega-\sin\theta$ con $\theta$ un ángulo tiene soluciones periódicas si $\omega>1$: la partícula da vueltas a la circunferencia.
  \item \emph{Unicidad.} Sin la condición de Lipschitz una solución puede detenerse en un equilibrio en tiempo finito, como $\dot x=-x^{1/3}$ (\cref{ej:un-raiz}).
\end{itemize}
\correcto{``No hay oscilaciones'' vale para $\dot x=f(x)$, $x\in\R$, con $f$ localmente Lipschitz. Las tres hipótesis cuentan.}
\end{trampa}

\subsection{La bolita que baja: el potencial}
\label{sec:un-potencial}

Hay otra manera de ver la \cref{prop:un-monotonia}, que anticipa el método directo de Lyapunov. Defina el \clave{potencial}
\begin{equation}
  V(x)\defeq-\int_0^xf(s)\,\dif s,\qquad\text{de modo que}\qquad\dot x=f(x)=-V'(x) .
  \label{eq:un-potencial}
\end{equation}
A lo largo de una solución,
\[
  \frac{\dif}{\dif t}V\bigl(x(t)\bigr)=V'(x)\,\dot x=-f(x)^2\le0 ,
\]
con igualdad solo en los equilibrios. La partícula siempre ``baja'' por la gráfica de $V$, como una bolita en un medio muy viscoso: los mínimos locales de $V$ son los equilibrios estables y los máximos, los inestables. Para $\dot x=\sin x$, $V(x)=\cos x-1$ (\cref{fig:un-potencial}). En el capítulo~7 el papel de $V$ lo hará una \clave{función de Lyapunov}: una función que decrece a lo largo de las soluciones aunque el sistema no sea un gradiente.

\begin{figure}[ht]
\centering
\includegraphics[width=\linewidth]{cap02/f_potencial}
\caption{Izquierda: potencial $V(x)=\cos x-1$ de $\dot x=\sin x$; los mínimos ($x=\pm\pi$) son los equilibrios estables. Derecha: $V(x(t))$ a lo largo de ocho soluciones; nunca crece.}
\label{fig:un-potencial}
\end{figure}

\codigorepo{cap02_unidimensionales/C2_Ej3_Potencial.m}{potencial}{El potencial a lo largo de las soluciones}

\begin{pausa}
¿Todo sistema escalar $\dot x=f(x)$ con $f$ continua tiene un potencial? ¿Y todo sistema en $\R^2$?
\respuesta{En $\R$ sí: $V=-\int f$ siempre existe. En $\R^2$ no: $\dot x=-\nabla V$ exige que el campo sea un gradiente, y un campo con rotación (como $\dot x_1=-x_2$, $\dot x_2=x_1$) no lo es. Por eso en dimensión mayor hay que \emph{buscar} una función que decrezca: es el problema de construir funciones de Lyapunov (capítulos~7 y~8).}
\end{pausa}
